\documentclass[11pt]{article}
\usepackage{amsmath,amssymb}
\title{Harmonic remainders and retained heat-equation contacts}\author{}\date{}
\begin{document}\maketitle
\paragraph{Scope.}
Use the original flat relative SU(2) complex, with its complete boundary
and the same normalization as the complete shortest time contact packet calculation, the internal ward and heat regulator contacts calculation and the complete harmonic kinetic matching calculation. Let $H=d_1^*d_1$ on free
links, $P$ its physical projection and $H^+$ its inverse on $P$.
Gauge zero modes are quotiented, not supplied with artificial mass.
The harmonic matching of the complete harmonic kinetic matching calculation and the complete constant curvature coarse sector calculation fixes a scalar $c_R$ on the
18-dimensional constant-curvature boundary sector,
$|c_R|\le C\log(2+R)$. It supplies no form bound on arbitrary internal
fields. The identities below expose exactly the missing terms.

\paragraph{1. Original harmonic decomposition and the current it leaves.}
For retained boundary data $v$ write the physical field
$x=h(v)+(r,0)$, where $h(v)$ solves the original interior Maxwell
problem and $r$ has zero boundary. The classical energy splits
orthogonally, with its original Schur energy on $v$. For the COMPLETE
uncut extended tensor $\widehat K$ of the complete harmonic kinetic matching calculation, however,
\[
 x^*\widehat Kx=h^*\widehat Kh+2\langle r,j(v)\rangle
                         +r^*P\widehat K_{JJ}Pr,
 \qquad j(v)=P(\widehat Kh)_J.                         \tag{1}
\]
All fixed-only boundary terms belong to the first term, not to the
interior tensor. Let $k(v)=H^+j(v)$. Then the exact cross term is
$2\langle r,Hk(v)\rangle$. Harmonicity of $h$ does NOT imply $j=0$.
The original gaussian coarse kinetic interface calculation controls the declared constant-curvature kinetic piece on
retained subblock boundaries; it does not bound $k(v)$ or the last term.

For a time-dependent current $k_s$ and a physical test $f$,
\[
 \langle e^{-sH}f,Hk_s\rangle
 =-\partial_s\langle e^{-sH}f,k_s\rangle
                       +\langle e^{-sH}f,\partial_s k_s\rangle. \tag{2}
\]
Integration keeps the $s=0$ endpoint $\langle f,k_0\rangle$ and the
second term. This is the equation-of-motion contact for an actual cut
leg. It is not an uncut Ward identity applied at each proper time.
If a background source varies $H,P,k_s$, all their derivatives must
also be differentiated in (2); no fixed-projector simplification is
made for that source response.

\paragraph{2. An exact kinetic add-and-subtract with its full remainder.}
Use the completed transverse packet of the exact completed transverse shortest time packet calculation:
\[
 C=\int_0^\infty\left\{\tfrac12\operatorname{Tr}_P
             (e^{-sH}K_s^\perp)+\operatorname{Tr}
             (e^{-s\Delta_0}J_s^\perp)\right\}\,ds.             \tag{3}
\]
Here $K_s^\perp=PK(G_s^\perp,g_s)P$,
$G_s^\perp=H^+e^{-sH}P$, and the ghost marking is retained.
The longitudinal total-time subtraction of the exact completed transverse shortest time packet calculation has endpoints zero only
after the full coefficient is assembled. Its differentiated versions
remain part of any background-source calculation.

One may subtract the DECLARED uncut matched value $c_R$ algebraically:
\[
 B_s=H^+(K_s^\perp-c_RH)H^+,
 \qquad K_s^\perp=c_RH+HB_sH.                              \tag{4}
\]
This definition does not assert that $c_R$ is the kinetic coefficient
of every cut packet or of every internal mode. Rather it fixes the
subtraction by the known harmonic matching and stores ALL mismatch in
$B_s$. At every finite box,
\[
 \int_0^\infty\operatorname{Tr}_P(e^{-sH}HB_sH)\,ds
 =\operatorname{Tr}_P(HB_0)
       +\int_0^\infty\operatorname{Tr}_P(He^{-sH}B_s')\,ds       \tag{5}
\]
and, equivalently,
\[
 =\operatorname{Tr}_P(HB_0+B_0')
       +\int_0^\infty\operatorname{Tr}_P(e^{-sH}B_s'')\,ds.    \tag{6}
\]
The proof is two finite-dimensional integrations by parts. Heat decay
makes the upper endpoints zero, although $B_s$ itself need not decay.
The original packet is polynomial in the decaying covariances, so all
needed derivatives exist and are integrable at fixed size.
The kinetic contribution is $c_R\dim(P)/2$ in this particular flat
trace channel. It is stored as a scalar, not discarded. Neither this
scalar nor the endpoint tensors in (5)--(6) are separately a source
response: the marked ghost, source changes, and longitudinal endpoint
identities still have to be combined. In particular, replacing (5) by
zero would lose the short-line contact under investigation.

\paragraph{3. A concrete limitation of an energy-squared remainder bound.}
The full tensor of the complete shortest time contact packet calculation includes the free Haar term $2I/3$. When both
covariances tend to zero all other terms vanish. Hence, on the physical
space,
\[
 K_s^\perp\longrightarrow\tfrac23P,\qquad
 B_s\longrightarrow\tfrac23(H^+)^2-c_RH^+.                    \tag{7}
\]
The relative four-dimensional box has physical eigenvalues of order
$R^{-2}$. This follows directly from its product complex: take the
lowest nonzero sine frequency in every coordinate; the four edge
components have the common eigenvalue
$\lambda_R=16\sin^2(\pi/(2R))$, and the three coefficient combinations
orthogonal to the scalar-gradient combination are co-closed physical
modes. On them $d_1^*d_1=\Delta_1$.
Consequently, for sufficiently large $R$,
\[
 \|B_\infty\|\ge
 \left|\frac{2}{3\lambda_R^2}-\frac{c_R}{\lambda_R}\right|
 \ge cR^4.                                                   \tag{8}
\]
Thus a uniform-in-time $O(\log R)$ bound on this unweighted
energy-squared remainder is FALSE for the actual completed tensor.
This is not divergence of the full coefficient: the tensor is
heat-weighted in (3), its source susceptibility must be completed with
all partners, and cancellation may occur between times or markings.
It does rule out proving the desired estimate merely by attaching two
inverse classical energies to the internal remainder and taking its
operator norm.

\paragraph{4. The remaining precise interface.}
For the constant-flux source, the useful still-unproved quantity is the
second source derivative of the COMBINED right side of (3), or its
exact form (4)--(6), including marked ghost and projector/minimum
partners. A bound must control the heat-weighted transverse remainder
and its zero-time contacts jointly, in physical curl norm. Source Ward
identities and the 18-mode harmonic matching do not do this alone:
whole-block zero-boundary heat legs have zero projection onto those
18 modes by literal Stokes. Genuine retained subblock boundaries can
use the original gaussian coarse kinetic interface calculation, with its exact Schur determinant and joint covariance, but the
residual currents (1) and internal terms (4) remain.
No complete nonlinear, physical-scale or S94 assertion follows.
\end{document}
